% !TeX program = xelatex
% Все численные координаты рисунков и таблиц включены в этот файл.
\documentclass[11pt,a4paper]{article}
\usepackage{fontspec}
\usepackage[russian]{babel}
\IfFontExistsTF{Cambria}{\setmainfont{Cambria}}{\setmainfont{DejaVu Serif}}
\IfFontExistsTF{Calibri}{\setsansfont{Calibri}}{\setsansfont{DejaVu Sans}}
\usepackage[a4paper,left=21mm,right=19mm,top=22mm,bottom=20mm,headheight=15pt,headsep=5mm,footskip=10mm]{geometry}
\usepackage{amsmath,booktabs,array,graphicx,xcolor,fancyhdr,hyperref,pgfplots,caption,titlesec}
\pgfplotsset{compat=1.16}
\definecolor{ink}{HTML}{193644}\definecolor{teal}{HTML}{066E70}\definecolor{rust}{HTML}{9E4F39}
\definecolor{plum}{HTML}{73569A}\definecolor{ochre}{HTML}{916316}\definecolor{sky}{HTML}{34659E}
\definecolor{muted}{HTML}{4D626A}\definecolor{hair}{HTML}{CFD9DA}
\hypersetup{unicode=true,colorlinks=true,linkcolor=teal,urlcolor=teal,pdftitle={Ti–Al. Замыкание сольвуса на MTP},pdfsubject={Атомистические расчёты при нулевом давлении}}
\captionsetup{font=small,labelfont={bf,color=teal},labelsep=period,singlelinecheck=false}
\titleformat{\section}{\sffamily\bfseries\Large\color{ink}}{}{0pt}{}
\titleformat{\subsection}{\sffamily\bfseries\large\color{ink}}{}{0pt}{}
\titlespacing*{\section}{0pt}{4pt}{7pt}\titlespacing*{\subsection}{0pt}{9pt}{3pt}
\setlength{\parindent}{0pt}\setlength{\parskip}{4pt}\setlength{\emergencystretch}{2em}
\renewcommand{\arraystretch}{1.18}\raggedbottom\pagestyle{fancy}\fancyhf{}
\fancyhead[L]{\small\sffamily\color{muted}Ti--Al\quad/\quad Равновесие на ГПУ-решётке}
\fancyhead[R]{\small\sffamily\color{muted}Прямой MTP\quad/\quad \(p=0\)}
\fancyfoot[R]{\small\sffamily\color{muted}\thepage}
\renewcommand{\headrulewidth}{.25pt}
\begin{document}
\thispagestyle{plain}
{\sffamily\bfseries\fontsize{23}{26}\selectfont\color{ink}Где замыкается сольвус \(\alpha/\alpha_2\) в Ti--Al?\par}
{\large\sffamily\color{muted}Расчёт на MTP Qi при нулевом давлении\par}
\vspace{1mm}{\color{teal}\rule{\linewidth}{1.1pt}}
На MTP Qi область упорядочения D0$_{19}$ замыкается около \textbf{1244~К и 27,2~ат.\% Al}. Рабочий диапазон для объёмного сплава при заданном ограничении смещений составляет \textbf{1240--1247~К}. Он учитывает разброс независимых серий и способов экстраполяции. На 2304 атомах условный 95\%-й интервал составляет 1242,4--1245,0~К и 27,12--27,23~ат.\% Al. При общем составе параметр порядка меняется на 0,66, а теплота разупорядочения равна 27,9~мэВ/атом (26,9--29,0). Сохранение двух структурных пиков и рост барьера с размером поддерживают конгруэнтный переход первого рода.

\subsection*{Температура и область, охваченная прямым расчётом}
\begin{center}\small\renewcommand{\arraystretch}{1.08}\setlength{\tabcolsep}{3.3pt}\begin{tabular}{@{}rrrrrrcc@{}}\toprule
\(T\), К & \(N\) & \(x_\alpha\) & \(x_{\alpha_2}\) & \(\Delta x\), 95\% & \(\beta\Delta F\) & Ti & Al\\\midrule
1000 & 1400 & 14,98 & 22,44 & [7,32; 7,58] & 5,66 & П & ---\\
1100 & 768 & 17,94 & 23,01 & [4,92; 5,21] & 3,22 & П & В\\
1150 & 768 & 19,70 & 23,39 & [3,58; 3,80] & 2,61 & П & В\\
1200 & 1400 & 22,00 & 24,31 & [2,19; 2,41] & 3,53 & П & В\\
1225 & 1400 & 23,68 & 25,08 & [1,32; 1,51] & 3,05 & П & В\\
1237,50 & 288 & --- & --- & --- & --- & Н & Н\\
1240 & 1400 & --- & --- & --- & --- & Н & Н\\
1245 & 2304 & --- & --- & --- & --- & Н & Н\\
1247,50 & 1400 & --- & --- & --- & --- & Н & Н\\
1250 & 1400 & --- & --- & --- & --- & Н & Н\\
1273,15 & 1400 & --- & --- & --- & --- & Н & ---\\
\bottomrule\end{tabular}\end{center}
{\small Составы в ат.\% Al на Ti-богатой ветви приведены для наибольшей ячейки изотермы. Для \(\Delta x=x_{\alpha_2}-x_\alpha\) показан условный 95\%-й интервал в п.п. Обозначения ветвей следующие. П --- подтверждённое сосуществование, К --- слабый кандидат, Н --- пара не разрешена, В --- недостаточное покрытие полей. Прочерк означает отсутствие оценки.}\par
При 1225~К ширина Ti-богатой ветви остаётся 1,40~п.п. для 768 и 1400 атомов, а барьер увеличивается с 1,98 до 3,05~\(k_{\rm B}T\). Разрешение ветвей при 1240~К зависит от выбора разделителя. При \(q_c=0{,}47\) на 1400 атомах оценки ширины равны \(+0{,}30\) и \(-0{,}25\)~п.п., однако 95\%-й интервал максимума логарифма весов включает ноль.

На 2304 атомах перенос с 1245 на 1243~К прогнозирует \(+0{,}23\) и \(-0{,}21\)~п.п. на двух ветвях. Оба корня найдены в 53 из 64 блочных повторов. Их поля выходят за границы прямой сетки на 4,5 и 2,3~мэВ соответственно. При переносе на 1240~К эффективное число кадров порядка единицы.
Оценка Al-богатой ветви (\(x_\alpha/x_{\alpha_2}\), ат.\% Al) получена при 1225 К, \(N=768\) --- 30,76/29,89 (кандидат).
\begin{center}\begin{tikzpicture}\begin{axis}[width=.91\linewidth,height=51mm,axis x line*=bottom,axis y line*=left,ymajorgrids,y grid style={hair!70},tick label style={font=\sffamily\footnotesize},label style={font=\sffamily\small},legend columns=4,legend style={font=\sffamily\scriptsize,draw=none,column sep=3pt,at={(.5,1.01)},anchor=south},enlarge y limits={upper,value=.13},height=40mm,xmin=0,xmax=1,ymin=0,ymax=5.5,xtick={0,.2,.4,.6,.8,1},xlabel={Параметр порядка \(Q\)},ylabel={\(-\ln[P(Q)/P_{\max}]\)}]
\addplot[sky,line width=1pt,no marks,unbounded coords=jump] coordinates {(0.0025,2.42153) (0.0075,2.40270) (0.0125,2.36647) (0.0175,2.31541) (0.0225,2.25277) (0.0275,2.18201) (0.0325,2.10635) (0.0375,2.02860) (0.0425,1.95106) (0.0475,1.87552) (0.0525,1.80328) (0.0575,1.73528) (0.0625,1.67214) (0.0675,1.61425) (0.0725,1.56181) (0.0775,1.51489) (0.0825,1.47346) (0.0875,1.43742) (0.0925,1.40660) (0.0975,1.38078) (0.1025,1.35973) (0.1075,1.34319) (0.1125,1.33087) (0.1175,1.32246) (0.1225,1.31766) (0.1275,1.31615) (0.1325,1.31762) (0.1375,1.32173) (0.1425,1.32819) (0.1475,1.33670) (0.1525,1.34696) (0.1575,1.35873) (0.1625,1.37176) (0.1675,1.38587) (0.1725,1.40088) (0.1775,1.41665) (0.1825,1.43309) (0.1875,1.45013) (0.1925,1.46771) (0.1975,1.48580) (0.2025,1.50439) (0.2075,1.52346) (0.2125,1.54298) (0.2175,1.56292) (0.2225,1.58323) (0.2275,1.60387) (0.2325,1.62474) (0.2375,1.64575) (0.2425,1.66681) (0.2475,1.68780) (0.2525,1.70861) (0.2575,1.72914) (0.2625,1.74930) (0.2675,1.76899) (0.2725,1.78816) (0.2775,1.80676) (0.2825,1.82479) (0.2875,1.84223) (0.2925,1.85912) (0.2975,1.87550) (0.3025,1.89140) (0.3075,1.90689) (0.3125,1.92203) (0.3175,1.93688) (0.3225,1.95148) (0.3275,1.96588) (0.3325,1.98010) (0.3375,1.99419) (0.3425,2.00810) (0.3475,2.02183) (0.3525,2.03534) (0.3575,2.04856) (0.3625,2.06144) (0.3675,2.07387) (0.3725,2.08576) (0.3775,2.09702) (0.3825,2.10756) (0.3875,2.11728) (0.3925,2.12608) (0.3975,2.13391) (0.4025,2.14069) (0.4075,2.14639) (0.4125,2.15097) (0.4175,2.15444) (0.4225,2.15676) (0.4275,2.15794) (0.4325,2.15795) (0.4375,2.15681) (0.4425,2.15449) (0.4475,2.15096) (0.4525,2.14617) (0.4575,2.14006) (0.4625,2.13259) (0.4675,2.12365) (0.4725,2.11320) (0.4775,2.10116) (0.4825,2.08751) (0.4875,2.07219) (0.4925,2.05523) (0.4975,2.03664) (0.5025,2.01650) (0.5075,1.99487) (0.5125,1.97185) (0.5175,1.94755) (0.5225,1.92212) (0.5275,1.89567) (0.5325,1.86832) (0.5375,1.84017) (0.5425,1.81134) (0.5475,1.78193) (0.5525,1.75197) (0.5575,1.72153) (0.5625,1.69064) (0.5675,1.65934) (0.5725,1.62766) (0.5775,1.59559) (0.5825,1.56317) (0.5875,1.53042) (0.5925,1.49731) (0.5975,1.46387) (0.6025,1.43010) (0.6075,1.39601) (0.6125,1.36158) (0.6175,1.32678) (0.6225,1.29160) (0.6275,1.25600) (0.6325,1.21994) (0.6375,1.18335) (0.6425,1.14617) (0.6475,1.10834) (0.6525,1.06983) (0.6575,1.03058) (0.6625,0.99058) (0.6675,0.94982) (0.6725,0.90835) (0.6775,0.86621) (0.6825,0.82351) (0.6875,0.78036) (0.6925,0.73691) (0.6975,0.69334) (0.7025,0.64984) (0.7075,0.60660) (0.7125,0.56384) (0.7175,0.52174) (0.7225,0.48051) (0.7275,0.44032) (0.7325,0.40133) (0.7375,0.36369) (0.7425,0.32750) (0.7475,0.29288) (0.7525,0.25989) (0.7575,0.22861) (0.7625,0.19907) (0.7675,0.17132) (0.7725,0.14539) (0.7775,0.12131) (0.7825,0.09912) (0.7875,0.07887) (0.7925,0.06065) (0.7975,0.04452) (0.8025,0.03061) (0.8075,0.01904) (0.8125,0.01000) (0.8175,0.00366) (0.8225,0.00025) (0.8275,-0.00000) (0.8325,0.00318) (0.8375,0.01006) (0.8425,0.02093) (0.8475,0.03610) (0.8525,0.05587) (0.8575,0.08055) (0.8625,0.11046) (0.8675,0.14588) (0.8725,0.18712) (0.8775,0.23447) (0.8825,0.28821) (0.8875,0.34862) (0.8925,0.41593) (0.8975,0.49041) (0.9025,0.57229) (0.9075,0.66178) (0.9125,0.75909) (0.9175,0.86444) (0.9225,0.97799) (0.9275,1.09989) (0.9325,1.23028) (0.9375,1.36923) (0.9425,1.51679) (0.9475,1.67283) (0.9525,1.83706) (0.9575,2.00891) (0.9625,2.18728) (0.9675,2.37030) (0.9725,2.55464) (0.9775,2.73502) (0.9825,2.90356) (0.9875,3.04907) (0.9925,3.15791) (0.9975,3.21651)};
\addlegendentry{\(N=768\)}
\addplot[teal,line width=1pt,no marks,unbounded coords=jump] coordinates {(0.0025,1.19497) (0.0075,1.18082) (0.0125,1.15359) (0.0175,1.11521) (0.0225,1.06816) (0.0275,1.01514) (0.0325,0.95876) (0.0375,0.90138) (0.0425,0.84496) (0.0475,0.79107) (0.0525,0.74090) (0.0575,0.69530) (0.0625,0.65486) (0.0675,0.61993) (0.0725,0.59070) (0.0775,0.56720) (0.0825,0.54937) (0.0875,0.53707) (0.0925,0.53011) (0.0975,0.52826) (0.1025,0.53123) (0.1075,0.53874) (0.1125,0.55049) (0.1175,0.56616) (0.1225,0.58546) (0.1275,0.60807) (0.1325,0.63367) (0.1375,0.66198) (0.1425,0.69270) (0.1475,0.72556) (0.1525,0.76029) (0.1575,0.79666) (0.1625,0.83442) (0.1675,0.87338) (0.1725,0.91336) (0.1775,0.95419) (0.1825,0.99575) (0.1875,1.03791) (0.1925,1.08060) (0.1975,1.12374) (0.2025,1.16728) (0.2075,1.21117) (0.2125,1.25537) (0.2175,1.29981) (0.2225,1.34446) (0.2275,1.38923) (0.2325,1.43402) (0.2375,1.47873) (0.2425,1.52320) (0.2475,1.56728) (0.2525,1.61077) (0.2575,1.65350) (0.2625,1.69526) (0.2675,1.73587) (0.2725,1.77512) (0.2775,1.81288) (0.2825,1.84901) (0.2875,1.88345) (0.2925,1.91613) (0.2975,1.94705) (0.3025,1.97624) (0.3075,2.00377) (0.3125,2.02974) (0.3175,2.05423) (0.3225,2.07738) (0.3275,2.09930) (0.3325,2.12006) (0.3375,2.13977) (0.3425,2.15847) (0.3475,2.17621) (0.3525,2.19298) (0.3575,2.20877) (0.3625,2.22353) (0.3675,2.23722) (0.3725,2.24979) (0.3775,2.26117) (0.3825,2.27133) (0.3875,2.28022) (0.3925,2.28785) (0.3975,2.29423) (0.4025,2.29939) (0.4075,2.30340) (0.4125,2.30631) (0.4175,2.30819) (0.4225,2.30913) (0.4275,2.30922) (0.4325,2.30849) (0.4375,2.30700) (0.4425,2.30477) (0.4475,2.30182) (0.4525,2.29811) (0.4575,2.29362) (0.4625,2.28829) (0.4675,2.28205) (0.4725,2.27482) (0.4775,2.26650) (0.4825,2.25699) (0.4875,2.24624) (0.4925,2.23414) (0.4975,2.22062) (0.5025,2.20565) (0.5075,2.18920) (0.5125,2.17125) (0.5175,2.15183) (0.5225,2.13097) (0.5275,2.10873) (0.5325,2.08519) (0.5375,2.06041) (0.5425,2.03449) (0.5475,2.00749) (0.5525,1.97945) (0.5575,1.95043) (0.5625,1.92042) (0.5675,1.88942) (0.5725,1.85738) (0.5775,1.82424) (0.5825,1.78986) (0.5875,1.75419) (0.5925,1.71706) (0.5975,1.67838) (0.6025,1.63803) (0.6075,1.59592) (0.6125,1.55201) (0.6175,1.50624) (0.6225,1.45864) (0.6275,1.40928) (0.6325,1.35826) (0.6375,1.30574) (0.6425,1.25186) (0.6475,1.19689) (0.6525,1.14103) (0.6575,1.08453) (0.6625,1.02768) (0.6675,0.97072) (0.6725,0.91392) (0.6775,0.85752) (0.6825,0.80174) (0.6875,0.74681) (0.6925,0.69292) (0.6975,0.64024) (0.7025,0.58892) (0.7075,0.53912) (0.7125,0.49095) (0.7175,0.44452) (0.7225,0.39991) (0.7275,0.35722) (0.7325,0.31650) (0.7375,0.27782) (0.7425,0.24125) (0.7475,0.20683) (0.7525,0.17464) (0.7575,0.14475) (0.7625,0.11725) (0.7675,0.09223) (0.7725,0.06983) (0.7775,0.05019) (0.7825,0.03348) (0.7875,0.01988) (0.7925,0.00961) (0.7975,0.00290) (0.8025,-0.00000) (0.8075,0.00117) (0.8125,0.00667) (0.8175,0.01678) (0.8225,0.03178) (0.8275,0.05194) (0.8325,0.07753) (0.8375,0.10880) (0.8425,0.14603) (0.8475,0.18946) (0.8525,0.23933) (0.8575,0.29587) (0.8625,0.35932) (0.8675,0.42988) (0.8725,0.50778) (0.8775,0.59321) (0.8825,0.68639) (0.8875,0.78751) (0.8925,0.89676) (0.8975,1.01433) (0.9025,1.14040) (0.9075,1.27516) (0.9125,1.41876) (0.9175,1.57139) (0.9225,1.73320) (0.9275,1.90435) (0.9325,2.08497) (0.9375,2.27516) (0.9425,2.47506) (0.9475,2.68468) (0.9525,2.90390) (0.9575,3.13246) (0.9625,3.36967) (0.9675,3.61424) (0.9725,3.86318) (0.9775,4.11148) (0.9825,4.34998) (0.9875,4.56276) (0.9925,4.72800) (0.9975,4.81934)};
\addlegendentry{\(N=1400\)}
\addplot[plum,line width=1pt,no marks,unbounded coords=jump] coordinates {(0.0025,1.14092) (0.0075,1.12890) (0.0125,1.10577) (0.0175,1.07318) (0.0225,1.03329) (0.0275,0.98846) (0.0325,0.94103) (0.0375,0.89317) (0.0425,0.84669) (0.0475,0.80311) (0.0525,0.76360) (0.0575,0.72905) (0.0625,0.70008) (0.0675,0.67711) (0.0725,0.66040) (0.0775,0.65007) (0.0825,0.64614) (0.0875,0.64855) (0.0925,0.65717) (0.0975,0.67183) (0.1025,0.69233) (0.1075,0.71841) (0.1125,0.74980) (0.1175,0.78621) (0.1225,0.82733) (0.1275,0.87283) (0.1325,0.92236) (0.1375,0.97557) (0.1425,1.03212) (0.1475,1.09164) (0.1525,1.15377) (0.1575,1.21817) (0.1625,1.28450) (0.1675,1.35245) (0.1725,1.42169) (0.1775,1.49196) (0.1825,1.56298) (0.1875,1.63451) (0.1925,1.70633) (0.1975,1.77823) (0.2025,1.85002) (0.2075,1.92151) (0.2125,1.99255) (0.2175,2.06297) (0.2225,2.13260) (0.2275,2.20131) (0.2325,2.26892) (0.2375,2.33532) (0.2425,2.40034) (0.2475,2.46384) (0.2525,2.52571) (0.2575,2.58578) (0.2625,2.64392) (0.2675,2.70001) (0.2725,2.75395) (0.2775,2.80560) (0.2825,2.85489) (0.2875,2.90171) (0.2925,2.94604) (0.2975,2.98782) (0.3025,3.02704) (0.3075,3.06373) (0.3125,3.09794) (0.3175,3.12972) (0.3225,3.15920) (0.3275,3.18651) (0.3325,3.21176) (0.3375,3.23513) (0.3425,3.25680) (0.3475,3.27693) (0.3525,3.29572) (0.3575,3.31330) (0.3625,3.32986) (0.3675,3.34554) (0.3725,3.36045) (0.3775,3.37470) (0.3825,3.38836) (0.3875,3.40146) (0.3925,3.41404) (0.3975,3.42605) (0.4025,3.43744) (0.4075,3.44810) (0.4125,3.45790) (0.4175,3.46665) (0.4225,3.47414) (0.4275,3.48012) (0.4325,3.48432) (0.4375,3.48643) (0.4425,3.48615) (0.4475,3.48320) (0.4525,3.47733) (0.4575,3.46830) (0.4625,3.45597) (0.4675,3.44025) (0.4725,3.42115) (0.4775,3.39874) (0.4825,3.37314) (0.4875,3.34459) (0.4925,3.31333) (0.4975,3.27959) (0.5025,3.24367) (0.5075,3.20583) (0.5125,3.16624) (0.5175,3.12512) (0.5225,3.08262) (0.5275,3.03876) (0.5325,2.99362) (0.5375,2.94714) (0.5425,2.89925) (0.5475,2.84986) (0.5525,2.79882) (0.5575,2.74595) (0.5625,2.69113) (0.5675,2.63413) (0.5725,2.57485) (0.5775,2.51317) (0.5825,2.44896) (0.5875,2.38224) (0.5925,2.31303) (0.5975,2.24142) (0.6025,2.16760) (0.6075,2.09179) (0.6125,2.01428) (0.6175,1.93540) (0.6225,1.85552) (0.6275,1.77501) (0.6325,1.69425) (0.6375,1.61359) (0.6425,1.53338) (0.6475,1.45392) (0.6525,1.37545) (0.6575,1.29823) (0.6625,1.22240) (0.6675,1.14811) (0.6725,1.07545) (0.6775,1.00450) (0.6825,0.93529) (0.6875,0.86785) (0.6925,0.80218) (0.6975,0.73828) (0.7025,0.67617) (0.7075,0.61585) (0.7125,0.55735) (0.7175,0.50074) (0.7225,0.44608) (0.7275,0.39349) (0.7325,0.34312) (0.7375,0.29513) (0.7425,0.24975) (0.7475,0.20721) (0.7525,0.16778) (0.7575,0.13174) (0.7625,0.09940) (0.7675,0.07106) (0.7725,0.04703) (0.7775,0.02762) (0.7825,0.01312) (0.7875,0.00383) (0.7925,-0.00000) (0.7975,0.00189) (0.8025,0.00974) (0.8075,0.02375) (0.8125,0.04413) (0.8175,0.07106) (0.8225,0.10470) (0.8275,0.14522) (0.8325,0.19275) (0.8375,0.24743) (0.8425,0.30941) (0.8475,0.37880) (0.8525,0.45574) (0.8575,0.54036) (0.8625,0.63277) (0.8675,0.73312) (0.8725,0.84154) (0.8775,0.95816) (0.8825,1.08312) (0.8875,1.21655) (0.8925,1.35862) (0.8975,1.50945) (0.9025,1.66919) (0.9075,1.83800) (0.9125,2.01601) (0.9175,2.20339) (0.9225,2.40028) (0.9275,2.60681) (0.9325,2.82313) (0.9375,3.04943) (0.9425,3.28576) (0.9475,3.53230) (0.9525,3.78904) (0.9575,4.05584) (0.9625,4.33255) (0.9675,4.61805) (0.9725,4.90982) (0.9775,5.20420) (0.9825,5.49056) (0.9875,5.75233) (0.9925,5.96129) (0.9975,6.08001)};
\addlegendentry{\(N=2304\)}
\end{axis}\end{tikzpicture}\end{center}
\captionof{figure}{Профили при 1245~К в поле максимума отношения весов при \(q_c=0{,}47\). Нуль выбран для наиболее вероятного \(Q\) каждого размера. Локальное эффективное число кадров после поправки на корреляции не ниже пяти.}\label{fig:order}
\clearpage
\section*{Составы и барьер на конечных размерах}
\begin{center}\begin{tikzpicture}\begin{axis}[width=.91\linewidth,height=51mm,axis x line*=bottom,axis y line*=left,ymajorgrids,y grid style={hair!70},tick label style={font=\sffamily\footnotesize},label style={font=\sffamily\small},legend columns=4,legend style={font=\sffamily\scriptsize,draw=none,column sep=3pt,at={(.5,1.01)},anchor=south},enlarge y limits={upper,value=.13},ymin=-1.5,ytick={-1,0,2,4,6,8},xmin=983,xmax=1286,xtick={1000,1050,1100,1150,1200,1250,1273.15},xticklabels={1000,1050,1100,1150,1200,1250,1273},xlabel={Температура, К},ylabel={\(\Delta x\), п.п. Al}]
\addplot[hair,dashed,no marks,forget plot] coordinates {(983,0) (1286,0)};
\addplot[no marks,forget plot,line width=.45pt,muted] coordinates {(1100.000000,5.215981) (1100.000000,5.509594)};
\addplot[no marks,forget plot,line width=.45pt,muted] coordinates {(1150.000000,3.859888) (1150.000000,4.394773)};
\addplot[no marks,forget plot,line width=.45pt,muted] coordinates {(1200.000000,2.555655) (1200.000000,3.125009)};
\addplot[no marks,forget plot,line width=.45pt,muted] coordinates {(1000.000000,6.987444) (1000.000000,7.722675)};
\addplot[only marks,forget plot,mark=*,mark size=2.5pt,muted,line width=.8pt] coordinates {(1000.000000,7.380136)};
\addplot[only marks,forget plot,mark=o,mark size=2.5pt,muted,line width=.8pt] coordinates {(1100.000000,5.397431) (1150.000000,4.156609) (1200.000000,2.735093)};
\addlegendimage{only marks,mark=*,mark size=1.7pt,muted}
\addlegendentry{\(N=96\)}
\addplot[no marks,forget plot,line width=.45pt,ochre] coordinates {(1100.000000,4.805088) (1100.000000,5.110183)};
\addplot[no marks,forget plot,line width=.45pt,ochre] coordinates {(1150.000000,3.652139) (1150.000000,3.919730)};
\addplot[no marks,forget plot,line width=.45pt,ochre] coordinates {(1200.000000,2.111923) (1200.000000,2.377608)};
\addplot[no marks,forget plot,line width=.45pt,ochre] coordinates {(1000.000000,6.954743) (1000.000000,7.449611)};
\addplot[only marks,forget plot,mark=triangle*,mark size=2.5pt,ochre,line width=.8pt] coordinates {(1000.000000,7.222496) (1100.000000,5.010423) (1150.000000,3.802139) (1200.000000,2.221104)};
\addlegendimage{only marks,mark=triangle*,mark size=1.7pt,ochre}
\addlegendentry{\(N=288\)}
\addplot[no marks,forget plot,line width=.45pt,sky] coordinates {(1100.000000,4.915980) (1100.000000,5.212648)};
\addplot[no marks,forget plot,line width=.45pt,sky] coordinates {(1150.000000,3.583442) (1150.000000,3.799933)};
\addplot[no marks,forget plot,line width=.45pt,sky] coordinates {(1200.000000,2.123124) (1200.000000,2.336590)};
\addplot[no marks,forget plot,line width=.45pt,sky] coordinates {(1225.000000,-0.917786) (1225.000000,-0.801196)};
\addplot[no marks,forget plot,line width=.45pt,sky] coordinates {(1225.000000,1.303799) (1225.000000,1.491640)};
\addplot[no marks,forget plot,line width=.45pt,sky] coordinates {(1000.000000,7.345078) (1000.000000,7.582930)};
\addplot[only marks,forget plot,mark=square*,mark size=2.5pt,sky,line width=.8pt] coordinates {(1000.000000,7.463209) (1100.000000,5.067806) (1150.000000,3.691949) (1200.000000,2.249728) (1225.000000,1.403764)};
\addplot[only marks,forget plot,mark=square,mark size=2.5pt,sky,line width=.8pt] coordinates {(1225.000000,-0.870702)};
\addlegendimage{only marks,mark=square*,mark size=1.7pt,sky}
\addlegendentry{\(N=768\)}
\addplot[no marks,forget plot,line width=.45pt,teal] coordinates {(1200.000000,2.194722) (1200.000000,2.407514)};
\addplot[no marks,forget plot,line width=.45pt,teal] coordinates {(1225.000000,1.318349) (1225.000000,1.507424)};
\addplot[no marks,forget plot,line width=.45pt,teal] coordinates {(1000.000000,7.323101) (1000.000000,7.579613)};
\addplot[only marks,forget plot,mark=diamond*,mark size=2.5pt,teal,line width=.8pt] coordinates {(1000.000000,7.463087) (1200.000000,2.314839) (1225.000000,1.400374)};
\addlegendimage{only marks,mark=diamond*,mark size=1.7pt,teal}
\addlegendentry{\(N=1400\)}
\end{axis}\end{tikzpicture}\end{center}
\captionof{figure}{Разность составов \(x_{\alpha_2}-x_\alpha\). Ti-богатой ветви соответствует положительный знак, Al-богатой --- отрицательный. Цвет обозначает размер ячейки. Закрашенный символ отмечает разрешённое сосуществование, полый --- слабый кандидат. Отрезки показывают условные блочные 95\%-е интервалы там, где они оценены. Пропуски оставлены пустыми.}
\vspace{1mm}
\begin{center}\begin{tikzpicture}\begin{axis}[width=.91\linewidth,height=51mm,axis x line*=bottom,axis y line*=left,ymajorgrids,y grid style={hair!70},tick label style={font=\sffamily\footnotesize},label style={font=\sffamily\small},legend columns=4,legend style={font=\sffamily\scriptsize,draw=none,column sep=3pt,at={(.5,1.01)},anchor=south},enlarge y limits={upper,value=.13},ymin=0,xtick={20.965931,43.610894,83.863725,125.146495,174.443577},xticklabels={96,288,768,1400,2304},xlabel={Число атомов \(N\), шкала \(N^{2/3}\)},ylabel={\(\beta\Delta F\)}]
\addplot[no marks,forget plot,line width=.5pt,teal!65] coordinates {(20.965931,1.183160) (43.610894,2.168164)};
\addplot[no marks,forget plot,line width=.5pt,teal!65] coordinates {(43.610894,2.168164) (83.863725,4.230915)};
\addplot[no marks,forget plot,line width=.5pt,teal!65] coordinates {(83.863725,4.230915) (125.146495,5.661120)};
\addplot[only marks,forget plot,mark=*,mark size=2.5pt,teal,line width=.8pt] coordinates {(20.965931,1.183160) (43.610894,2.168164) (83.863725,4.230915) (125.146495,5.661120)};
\addlegendimage{only marks,mark=*,mark size=1.7pt,teal}
\addlegendentry{1000 К}
\addplot[no marks,forget plot,line width=.5pt,sky!65] coordinates {(20.965931,1.010807) (43.610894,1.633133)};
\addplot[no marks,forget plot,line width=.5pt,sky!65] coordinates {(43.610894,1.633133) (83.863725,3.216427)};
\addplot[no marks,forget plot,line width=.45pt,sky] coordinates {(20.965931,0.888028) (20.965931,1.091926)};
\addplot[no marks,forget plot,line width=.45pt,sky] coordinates {(43.610894,1.441085) (43.610894,1.804957)};
\addplot[no marks,forget plot,line width=.45pt,sky] coordinates {(83.863725,2.813158) (83.863725,3.599131)};
\addplot[only marks,forget plot,mark=triangle*,mark size=2.5pt,sky,line width=.8pt] coordinates {(43.610894,1.633133) (83.863725,3.216427)};
\addplot[only marks,forget plot,mark=triangle,mark size=2.5pt,sky,line width=.8pt] coordinates {(20.965931,1.010807)};
\addlegendimage{only marks,mark=triangle*,mark size=1.7pt,sky}
\addlegendentry{1100 К}
\addplot[no marks,forget plot,line width=.5pt,ochre!65] coordinates {(20.965931,0.910952) (43.610894,1.446981)};
\addplot[no marks,forget plot,line width=.5pt,ochre!65] coordinates {(43.610894,1.446981) (83.863725,2.613938)};
\addplot[no marks,forget plot,line width=.45pt,ochre] coordinates {(20.965931,0.777806) (20.965931,1.029526)};
\addplot[no marks,forget plot,line width=.45pt,ochre] coordinates {(43.610894,1.318859) (43.610894,1.624098)};
\addplot[no marks,forget plot,line width=.45pt,ochre] coordinates {(83.863725,2.352504) (83.863725,3.002455)};
\addplot[only marks,forget plot,mark=square*,mark size=2.5pt,ochre,line width=.8pt] coordinates {(43.610894,1.446981) (83.863725,2.613938)};
\addplot[only marks,forget plot,mark=square,mark size=2.5pt,ochre,line width=.8pt] coordinates {(20.965931,0.910952)};
\addlegendimage{only marks,mark=square*,mark size=1.7pt,ochre}
\addlegendentry{1150 К}
\addplot[no marks,forget plot,line width=.5pt,plum!65] coordinates {(20.965931,0.851052) (43.610894,1.292800)};
\addplot[no marks,forget plot,line width=.5pt,plum!65] coordinates {(43.610894,1.292800) (83.863725,2.331213)};
\addplot[no marks,forget plot,line width=.5pt,plum!65] coordinates {(83.863725,2.331213) (125.146495,3.525836)};
\addplot[no marks,forget plot,line width=.45pt,plum] coordinates {(20.965931,0.758945) (20.965931,0.985738)};
\addplot[no marks,forget plot,line width=.45pt,plum] coordinates {(43.610894,1.205605) (43.610894,1.415609)};
\addplot[no marks,forget plot,line width=.45pt,plum] coordinates {(83.863725,2.009565) (83.863725,2.666176)};
\addplot[no marks,forget plot,line width=.45pt,plum] coordinates {(125.146495,3.061514) (125.146495,4.170413)};
\addplot[only marks,forget plot,mark=diamond*,mark size=2.5pt,plum,line width=.8pt] coordinates {(43.610894,1.292800) (83.863725,2.331213) (125.146495,3.525836)};
\addplot[only marks,forget plot,mark=diamond,mark size=2.5pt,plum,line width=.8pt] coordinates {(20.965931,0.851052)};
\addlegendimage{only marks,mark=diamond*,mark size=1.7pt,plum}
\addlegendentry{1200 К}
\addplot[no marks,forget plot,line width=.5pt,rust!65] coordinates {(83.863725,1.982151) (125.146495,3.054921)};
\addplot[no marks,forget plot,line width=.45pt,rust] coordinates {(83.863725,1.762434) (83.863725,2.197796)};
\addplot[no marks,forget plot,line width=.45pt,rust] coordinates {(125.146495,2.621940) (125.146495,3.441630)};
\addplot[only marks,forget plot,mark=pentagon*,mark size=2.5pt,rust,line width=.8pt] coordinates {(83.863725,1.982151) (125.146495,3.054921)};
\addlegendimage{only marks,mark=pentagon*,mark size=1.7pt,rust}
\addlegendentry{1225 К}
\addplot[no marks,forget plot,line width=.45pt,black] coordinates {(43.610894,0.999345) (43.610894,1.283770)};
\addplot[only marks,forget plot,mark=x,mark size=2.5pt,black,line width=.8pt] coordinates {(43.610894,1.164429)};
\addlegendimage{only marks,mark=*,mark size=1.7pt,black}
\addlegendentry{1237,50 К}
\addplot[only marks,forget plot,mark=x,mark size=2.5pt,green!45!black,line width=.8pt] coordinates {(125.146495,2.208030) (174.443577,3.283466)};
\addlegendimage{only marks,mark=triangle*,mark size=1.7pt,green!45!black}
\addlegendentry{1245 К}
\addplot[only marks,forget plot,mark=x,mark size=2.5pt,muted,line width=.8pt] coordinates {(20.965931,0.778448) (43.610894,0.997506) (83.863725,1.550287)};
\addlegendimage{only marks,mark=square*,mark size=1.7pt,muted}
\addlegendentry{1250 К}
\addplot[only marks,forget plot,mark=x,mark size=2.5pt,blue!55!black,line width=.8pt] coordinates {(20.965931,0.637284)};
\addlegendimage{only marks,mark=diamond*,mark size=1.7pt,blue!55!black}
\addlegendentry{1273,15 К}
\end{axis}\end{tikzpicture}\end{center}
\captionof{figure}{Барьер на Ti-богатой стороне между равновысокими максимумами плотности по \(Q\). Цвет обозначает температуру. Линии соединяют соседние размеры, для которых оценены составы. Крестик обозначает пару максимумов \(Q\) без подтверждения отдельных составов при заданном отношении весов. Отсутствие оценки оставляет соответствующую позицию пустой.}\label{fig:eqbarrier}

При равных суммарных весах двух фаз химическая восприимчивость \(\partial\langle x\rangle/\partial\Delta\mu=\beta N\operatorname{Var}(x)\) содержит межфазный вклад \(\beta N(\Delta x)^2/4\), если \(\Delta x\) выражено в атомных долях. При сближении составов этот вклад исчезает. Продолжающийся в этой области разрыв порядка можно обнаружить по профилю \(Q\) и по энергии двух фаз.

При 1245~К в поле максимального отношения весов (\(q_c=0{,}47\)) барьер от меньшего пика растёт с 0,84 до 1,78 и 2,84~\(k_{\rm B}T\) для 768, 1400 и 2304 атомов. На 2304 атомах моды расположены при \(Q=0{,}083\) и 0,793 (рис.~\ref{fig:order}). При 1250~К на 768 атомах обе моды всё ещё наблюдаются, но максимальный вес одного варианта D0$_{19}$ составляет лишь 0,30 от веса разупорядоченного бассейна. Равновесная пара по отношению \(4:1\) в этой серии отсутствует.
\subsection*{Что остаётся проверить}
При 1273,15~К на 288 атомах сужение стенки \(W_{\max}\) с 0,40 до 0,35 при поле 3,48~эВ повысило \(\langle Q\rangle\) с 0,29 до 0,40 и состав на 0,31~п.п. Чтобы оценить вызванный стенкой сдвиг \(T_*\), нужен такой же контроль непосредственно у вершины. Более крупные ячейки с анализом геометрии границы помогут проверить макроскопическую скрытую теплоту и оценить межфазное натяжение.
\clearpage
\section*{Как определялись равновесные состояния}
Полубольшой изобарический Монте-Карло с \texttt{fitted.mtp} из набора Qi учитывает взаимные замены Ti и Al и классические ангармонические смещения. Относительные свободные энергии получены из равновесных вероятностей. В проходе выполняются по \(N\) проб замены и смещения, одна проба объёма, а раз в три прохода --- проба \(\log(c/a)\). В весе изменения объёма учитывается якобиан \((V'/V)^N\). Область ГПУ-кристалла задана условием \(W<0{,}40\). Здесь \(W=\max_{i,\mathbf d}2(\mathbf u_i-\bar{\mathbf u})\cdot\mathbf d/|\mathbf d|^2\), где \(\mathbf u_i\) --- смещение от узла, а \(\mathbf d\) пробегает двенадцать связей идеальной решётки. Допустимы объёмы 15--21~\AA$^3$/атом, \(1{,}45\le c/a\le1{,}72\); для проб формы и объёма требуется \(d_{\rm nn}^{\rm ideal}(1-W_{\max})\ge1{,}68052\)~\AA. Границы защищают от нефизического короткодействующего поведения MTP. \href{https://doi.org/10.1103/PhysRevMaterials.7.103602}{Qi и соавт. (2023)}, \href{https://doi.org/10.6084/m9.figshare.24079611}{потенциал и данные}.

В план расчётов входят ячейки из 96, 288, 768, 1400 и 2304 атомов. Соответствующие размеры в двухатомных ячейках ГПУ равны \(4{\times}4{\times}3\), \(6{\times}6{\times}4\), \(8{\times}8{\times}6\), \(10{\times}10{\times}7\) и \(12{\times}12{\times}8\). Через \(n_j\) обозначим число атомов Al на подрешётке \(j\), через \(k=\sum_jn_j\) их полное число. В каждое окно добавлен сдвиг энергии всей ячейки
\[
 U_{\rm bias}=\tfrac12 K(Q-q_0)^2-\Delta\mu\, k,\qquad
 Q=\frac{4(4\max_jn_j-k)}{3N}.
\]
Здесь \(\Delta\mu\) --- конфигурационное химическое поле на одну замену Ti на Al. Метод MBAR объединяет окна одного размера, температуры и механического ансамбля. Для проверки выборок учитываются временные корреляции, перекрытие соседних окон и вклад обеих фаз в оценку весов. \href{https://doi.org/10.1063/1.2978177}{Shirts и Chodera (2008).}

Поле \(\Delta\mu^*\) для составов выбирается по весам фаз. Четыре эквивалентных варианта D0$_{19}$ объединены операцией максимума в \(Q\), поэтому применяется условие \(P(\alpha_2)/P(\alpha)=4\). По разные стороны минимума между модами усредняется \(x=k/N\). Для оценки барьера поле отдельно подстраивается до равенства высот двух пиков плотности \(P(Q)\). Тогда \(\beta\Delta F=\ln(P_{\max}/P_{\min})\). Если седловая область между хорошо разделёнными фазами содержит две плоские границы при периодических условиях, главный вклад в барьер пропорционален их площади, порядка \(N^{2/3}\). Доменная геометрия, когерентные напряжения и корреляция \(Q\) со скрытыми степенями свободы могут изменить зависимость на малых ячейках. \href{https://doi.org/10.1103/PhysRevLett.65.137}{Lee и Kosterlitz (1990)}; \href{https://doi.org/10.1103/PhysRevB.47.14757}{Janke (1993).}

\subsection*{Конгруэнтное разупорядочение}
При фиксированном разделителе \(q_c\) обозначим \(R=P(Q>q_c)/[4P(Q\leq q_c)]\). Вершину области упорядочения ищем при температуре, где внутренний максимум \(\ln R\) по химическому полю равен нулю. На Ti-богатой ветви \(\ln R\) пересекает ноль снизу вверх с ростом поля, на Al-богатой --- сверху вниз. В точке слияния корней
\[
 \frac{\partial\ln R}{\partial\Delta\mu}=\beta N\,
      (x_{\alpha_2}-x_\alpha)=0.
\]
Совпадение составов в максимуме следует из этого уравнения. Независимые признаки разрыва --- разделённые моды порядка и разность энергий двух бассейнов. При одинаковом \(x\) и \(p=0\) оценка теплоты разупорядочения на атом принимает вид
\[
  \ell_{\rm dis} =
   \frac{\langle E-\Delta\mu\, k\rangle_\alpha -
         \langle E-\Delta\mu\, k\rangle_{\alpha_2}}{N}.
\]
Слагаемое \(\Delta\mu\,\Delta x\) исчезает, классические кинетические вклады фаз одинаковы. Положительное \(\ell_{\rm dis}\) соответствует поглощению теплоты. Модель концентрационных волн для ГПУ Ti--Al также допускает скачок порядка при стехиометрии. \href{https://arxiv.org/pdf/1406.0162}{Radchenko и соавт., рис.~4a.}

\subsection*{Переход к объёмному пределу}
Подгонка \(T_*(N)=T_\infty+a/N\) по 768, 1400 и 2304 атомам даёт 1242,5~К [1240,4; 1244,7]. Замена \(1/N\) на \(N^{-2/3}\) даёт 1242,4~К [1239,4; 1245,5]. По двум наибольшим ячейкам зависимость от \(1/N\) даёт 1246,5~К [1242,5; 1251,0]. Указанные 95\%-е блочные интервалы условны по форме каждой подгонки. Рабочая полоса 1240--1247~К характеризует межсерийный разброс и центральные экстраполяции; её статистическое покрытие неизвестно.
\clearpage
\section*{Независимые ориентиры и пределы вывода}
\subsection*{Проверка переноса по температуре}
При \(q_c=0{,}37\) ноль \(\max_{\Delta\mu}\ln R\) найден совместным перевзвешиванием по энергии и составу. После снятия зонтичного потенциала при \((\beta_0,\Delta\mu_0)\) логарифм поправки к весу кадра равен \(-(\beta-\beta_0)E+(\beta\Delta\mu-\beta_0\Delta\mu_0)k\). В совместном расчёте MBAR нормировки всех зонтиков разных температур оцениваются по исходным кадрам. Блочные интервалы условны по размеру, траекториям и разделителю.
\begin{center}\small\setlength{\tabcolsep}{3.6pt}\begin{tabular}{@{}rrrrrrr@{}}\toprule
\(T_0\), К & \(N\) & \(T_*\), К & 95\% для \(T_*\), К & \(x_*\) & \(\ell_{\rm dis}\) & \(B_{\rm low}\)\\\midrule
1250 & 288 & 1247,9 & [1244,7; 1251,3] & 27,30 & 29,8 & 0,34\\
1245 & 768 & 1244,4 & [1242,5; 1246,9] & 27,18 & 27,1 & 0,81\\
1245 & 1400 & 1242,3 & [1240,5; 1244,1] & 27,18 & 27,8 & 1,56\\
1245 & 2304 & 1243,6 & [1242,4; 1245,0] & 27,18 & 27,9 & 2,64\\
\(\sum T\) & 768 & 1243,8 & [1242,2; 1245,6] & 27,18 & 27,5 & 1,04\\
\(\sum T\) & 1400 & 1241,8 & [1240,4; 1243,3] & 27,17 & 27,8 & 1,57\\
\bottomrule\end{tabular}\end{center}
{\small \(T_0\) --- температура прямых траекторий, \(\sum T\) обозначает совместный MBAR по изотермам, \(x_*\) в ат.\% Al, \(\ell_{\rm dis}\) в мэВ/атом. \(B_{\rm low}\) --- глубина минимума в единицах \(k_{\rm B}T\) от меньшего \(Q\)-пика при \(R=1\). На рис.~\ref{fig:eqbarrier} поле отдельно подстроено до равных высот.}\par
На 2304 атомах первая и вторая половины траекторий дают 1242,7 и 1244,6~К. Изменение \(q_c\) с 0,32 до 0,47 сдвигает \(T_*\) на 0,17~К; на 288 атомах сдвиг около 7~К. На 768 атомах независимые переносы с 1225, 1250 и 1273,15~К дают 1246,6, 1243,1 и 1247,0~К.

На 2304 атомах локальный температурный наклон \(\max_{\Delta\mu}\ln R\) при 1245~К и \(q_c=0{,}47\) даёт 1243,52~К [1241,83; 1245,18]. Он проверяет нелинейное перевзвешивание на тех же траекториях другим способом. В совместном MBAR учёт каждого четвёртого вместо каждого восьмого кадра меняет \(T_*\) на 0,25--0,28~К.
\subsection*{Что измеряет гистерезис}
Nichols и соавт. наблюдали самопроизвольный переход на MTP Qi по траекториям длиной 40~пс на 1024 атомах. Составы перед переходами из \(\alpha\) и D0$_{19}$ равны 15,2/23,4~ат.\% Al при 1000~К и 18,2/23,7 при 1100~К. При 1200~К в табл.~2 работы стоят прочерки. Наши равновесные составы при 1000~К равны 14,98/22,44~ат.\% Al. Границы гистерезиса зависят также от времени наблюдения и начального состояния. \href{https://doi.org/10.3390/met16020140}{Nichols и соавт. (2026), разд.~2.4, табл.~2.}

\subsection*{Кластерное разложение и эксперимент}
При 1200~К квантовый C46 на том же MTP даёт 23,46/24,91~ат.\% Al (\(N=5120\)), прямой MTP --- 22,00/24,31~ат.\% Al (\(N=1400\)). Ширина области составляет 1,44 и 2,31~п.п. соответственно. В квантовом C46 при 1225, 1250 и 1300~К отдельные контакты не разрешены. В классическом C46 при 1273,15~К химический потенциал монотонен от 13,98 до 36,86~ат.\% Al. В сравнение входят усечение кластерного базиса, квазигармоническое приближение и фиксация \(c/a\) в C46. Квантовая статистика колебаний также отличается от классической, использованной при прямом расчёте.

Для диффузионной пары, отожжённой при 1273,15~К в течение 200~ч, электронно-зондовый микроанализ дал 19,5\(\pm\)0,5 и 23,0\(\pm\)0,3~ат.\% Al. При 1323,15~К получены 21,3\(\pm\)0,3 и 23,5\(\pm\)0,3, а при 1423,15~К --- 24,7\(\pm\)0,2 и 25,7\(\pm\)0,2~ат.\% Al. Это измерения фазовых границ на указанных изотермах. Знаки \(\pm\) передают точность определения составов. \href{https://doi.org/10.1016/0966-9795(94)90018-3}{Kainuma, Palm и Inden (1994), табл.~1.}

Для CALPHAD-модели Liang составы фаз совпадают при 1437,28~К и 29,76~ат.\% Al. Эта расчётная вершина выше оценки прямого MTP примерно на 193~К и богаче алюминием на 2,6~п.п. Измеренная пара при 1423,15~К непосредственно показывает сохранение двухфазной области на 179~К выше вершины MTP. \href{https://doi.org/10.1007/s11669-024-01124-9}{Liang и соавт. (2024), табл.~3.}

{\footnotesize\color{muted}Проверены 36 заданий UMT, 972 траектории и 738\,700 продуктивных кадров. Численные данные сохранены в \texttt{analysis/closure\_summary.csv}, \texttt{analysis/critical\_order\_profiles.csv} и \texttt{analysis/thermal/}. Координаты графиков встроены в исходник PDF.}
\end{document}
